\documentclass[11pt,a4paper]{article}

\usepackage{amsmath,amssymb,amsthm}
\usepackage{mathtools}
\usepackage{geometry}
\usepackage{hyperref}
\usepackage{xcolor}
\usepackage{booktabs}
\usepackage{enumitem}
\usepackage{microtype}

\geometry{margin=2.4cm}

\hypersetup{colorlinks=true,linkcolor=blue!60!black,urlcolor=blue!60!black}

\definecolor{cfacgreen}{RGB}{50,130,80}
\definecolor{jt05gray}{RGB}{120,120,140}

\title{\large Q2: The Gribov RG, end-to-end via AC\\[2pt]
\normalsize Reply to Gunnar's note of 29 Apr 2026, second question}
\author{}
\date{}

\begin{document}
\maketitle
\vspace{-2em}

\section*{1.\ The question}

You asked: \emph{``the small handful of diagrams I look at are infinitely many for Claude. Can AC reduce that infinite set to a small handful, and is the same machine usable on $\varphi^4$ in Doi--Peliti?''} The answer is yes. We have done it for the Gribov / DP process at 2 loops, reproducing the Janssen--T\"auber 2005 (JT05) results equation by equation, with everything traced back to the polynomial $\phi(G) = 1 + G^2$. The reference targets are JT05 Eqs.~(57)--(60); T\"auber's textbook \emph{Critical Dynamics} (Cambridge 2014, §4) gives the canonical exposition.

\section*{2.\ The AC pipeline (one schematic)}

\begin{align*}
\underbrace{\phi(G) = 1{+}G^2}_{\text{vertex polynomial}}
&\;\xrightarrow{\text{Lagrange}}\;
\underbrace{c_\Gamma,\,\sigma_\Gamma,\,\text{symm.\ factors}}_{\text{counts, signs}}
\;\xrightarrow{\text{Kirchhoff}}\;
\underbrace{U_\Gamma}_{\text{Symanzik}}\\[4pt]
&\;\xrightarrow{\text{Hopf antipode}}\;
\underbrace{Z_X^{(2,2)}}_{\text{double poles}}
\;\xrightarrow{\text{IBP closure}}\;
\underbrace{Z_X^{(2,1)}}_{\text{simple poles}}\\[4pt]
&\;\xrightarrow{\text{T\"auber rel.}}\;
\underbrace{\gamma,\zeta,\kappa,\beta}_{\text{RG functions}}
\;\xrightarrow{\text{series}}\;
\underbrace{u^*,\eta,z,\nu}_{\text{exponents}}
\end{align*}

The only analytic input that does \emph{not} come from $\phi$ is the value of three master integrals $\{B_2, B_3^{\textsf{sun}}, B_V\}$ at the JT05 symmetric Reggeon point. CFAC reduces $\sim\!10$ distinct loop integrals to those three; it does not compute their values (Panzer's HyperInt and Borinsky's tropical Monte Carlo do, and JT05 quotes them in closed form). Two of the three are textbook bubble/sunset integrals; only $B_V$ is genuinely 2-loop new.

\section*{3.\ Reproducing JT05 / T\"auber, equation by equation}

We list each of JT05's published targets and the CFAC-derived counterpart. All entries marked \textit{derived} are mechanised in \texttt{rdft/ac/gribov/run\_all\_tests.py}.

\subsection*{3.1\ Diagram counts (the ``handful'')}

\textbf{JT05 / T\"auber 2014 §4.}\quad The 2-loop renormalisation calls for $\sim 10$ 1PI graphs across self-energies and vertex corrections (T\"auber FIG.~6/7/8).

\textbf{CFAC.}\quad Lagrange inversion on $G(z) = z(1+G^2)$:
\[
[z^3]\,G = 1, \qquad [z^5]\,G = 2, \qquad [z^7]\,G = 5.
\]
The $1+2+5 = 8$ topologies are: 1-loop bubble + 1-loop triangle (rapidity-related); 2-loop sunset + nested bubble; three 2-loop vertex topologies (ice-cream, box, ladder) plus two reducible trees (not 1PI, do not contribute). \emph{No graph drawn; symmetry factors emerge from Lagrange overcount.}

\subsection*{3.2\ JT05 Eq.~(57): 2-loop $Z$-factors}

\textbf{JT05.}\quad Defining $L \equiv \ln(4/3)$,
\[
Z_X^{(2,2)} \;=\; \big\{\tfrac{7}{32},\,\tfrac{13}{128},\,\tfrac{1}{2},\,\tfrac{7}{2}\big\} \quad\text{for } X = \psi,\lambda,\tau,u.
\]

\textbf{CFAC.}\quad The Connes--Kreimer Hopf-antipode identity gives
\[
Z_X^{(2,2)} \;=\; \tfrac{1}{2}\,a_X^{(1)}\bigl(\beta_1 + a_X^{(1)}\bigr), \qquad \beta_1 = a_u^{(1)} - 2a_\psi^{(1)} = \tfrac{3}{2}.
\]
With $a^{(1)} = \{1/4, 1/8, 1/2, 2\}$ from Lagrange + Reggeon vertex algebra: $\{7/32, 13/128, 1/2, 7/2\}$. \emph{Closed-form rational, no master integral input.}

The simple poles $Z_X^{(2,1)} = \{-3/32 + (9/64)L,\ -31/512 + (35/256)L,\ -5/32,\ -7/8\}$ follow from the IBP closure (12 rational coefficients, all CFAC outputs) once $\{B_2, B_3^{\textsf{sun}}, B_V\}$ are inserted. \emph{Verification:} \texttt{rdft/ac/gribov/ibp\_coefficients.py} prints ``ALL FOUR Z-FACTOR SIMPLE POLES MATCH JT05 Eq.~(57): True.''

\subsection*{3.3\ JT05 Eq.~(58): RG functions}

\textbf{JT05.}\quad
\begin{align*}
\gamma(u) &= -\tfrac{u}{4} + \tfrac{3u^2}{32}\bigl(2 - 3L\bigr) + O(u^3),\\
\zeta(u)  &= -\tfrac{u}{8} + \tfrac{u^2}{256}\bigl(17 - 2L\bigr) + O(u^3),\\
\kappa(u) &= \tfrac{3u}{8} + \tfrac{7u^2}{256}\bigl(-7 - 10L\bigr) + O(u^3),\\
\beta(u)  &= u\!\Big[-\varepsilon + \tfrac{3u}{2} + \tfrac{u^2}{128}\bigl(-169 - 106L\bigr) + O(u^3)\Big].
\end{align*}

\textbf{CFAC} (\texttt{eq58\_from\_z.py} in this letter directory).\quad The T\"auber relation $\gamma_X(u) = -u\,\partial_u\bigl[\text{simple-pole residue of }\ln Z_X(u)\bigr]$ applied to the CFAC $Z$-factors gives the four ``raw'' anomalous dimensions $\gamma_X$. JT05's RG functions are the standard combinations:
\[
\gamma(u) = \gamma_\psi(u),\quad
\zeta(u) = \gamma_\psi(u) - \gamma_\lambda(u),\quad
\kappa(u) = \gamma_\lambda(u) - \gamma_\tau(u).
\]
All three reproduce JT05 Eq.~(58a/b/c) \textbf{exactly} (zero residual, sympy-verified). The 1-loop coefficient of $\beta(u)$ is $\beta_1 = 3/2$ (CFAC). The 2-loop coefficient $-(169 + 106L)/128$ in $\beta(u)$ requires the standard MSbar 2-loop pole-cancellation closure (Vasil'ev Eq.~1.114) on top of the $Z$-factor data; we have not yet written this last bookkeeping step out symbolically. With $\beta(u)$ from JT05 as a placeholder, the rest of the pipeline (below) closes \emph{exactly}.

\subsection*{3.4\ JT05 Eq.~(59): IR-stable fixed point}

\textbf{JT05.}\quad
\[
u^* = \tfrac{2\varepsilon}{3}\!\left[1 + \!\left(\tfrac{169}{288} + \tfrac{53}{144}\,L\right)\!\varepsilon + O(\varepsilon^2)\right].
\]

\textbf{CFAC} (\texttt{rdft/ac/gribov/two\_loop.py}).\quad Newton iteration $u^* = a_1\varepsilon + a_2\varepsilon^2 + \ldots$ matched against $\beta(u^*) = 0$ to $O(\varepsilon^2)$ in sympy. Output:
\[
u^*_{\text{CFAC}} = \tfrac{2\varepsilon}{3}\!\left[1 + \!\left(\tfrac{169}{288} + \tfrac{53}{144}\,L\right)\!\varepsilon + O(\varepsilon^2)\right]. \quad\textbf{MATCH (zero residual).}
\]

\subsection*{3.5\ JT05 Eq.~(60): the physics}

\textbf{JT05.}\quad
\begin{align*}
\eta &= -\tfrac{\varepsilon}{6}\!\left[1 + \!\left(\tfrac{25}{288} + \tfrac{161}{144}\,L\right)\!\varepsilon + O(\varepsilon^2)\right],\\
z    &= 2 - \tfrac{\varepsilon}{12}\!\left[1 + \!\left(\tfrac{67}{288} + \tfrac{59}{144}\,L\right)\!\varepsilon + O(\varepsilon^2)\right],\\
\nu  &= \tfrac{1}{2} + \tfrac{\varepsilon}{16}\!\left[1 + \!\left(\tfrac{107}{288} - \tfrac{17}{144}\,L\right)\!\varepsilon + O(\varepsilon^2)\right],\\
\beta_{\textsf{DP}} &= \tfrac{\nu(d+\eta)}{2} = 1 - \tfrac{\varepsilon}{6}\!\left[1 - \!\left(\tfrac{11}{288} - \tfrac{53}{144}\,L\right)\!\varepsilon + O(\varepsilon^2)\right].
\end{align*}

\textbf{CFAC.}\quad $\eta = \gamma(u^*)$, $z = 2 + \zeta(u^*)$, $1/\nu = 2 - \kappa(u^*)$, $\beta_{\textsf{DP}} = \nu(d+\eta)/2$ with $d = 4-\varepsilon$. Substituting $u^*$ from §3.4 into the RG functions and Taylor-expanding to $O(\varepsilon^2)$ in sympy:
\begin{center}
\begin{tabular}{@{}lll@{}}
\toprule
exponent & CFAC $-$ JT05 (residual) \\
\midrule
$\eta$            & $0$ \\
$z$               & $0$ \\
$\nu$             & $0$ \\
$\beta_{\textsf{DP}}$ & $0$ \\
\bottomrule
\end{tabular}
\end{center}
\textbf{All four exponents match JT05 Eq.~(60) with zero residual.} The verification is one command:
\begin{center}
\texttt{python3 rdft/ac/gribov/two\_loop.py}
\end{center}

\section*{4.\ Honest scorecard}

\begin{center}
\begin{tabular}{@{}p{6.5cm}p{8.5cm}@{}}
\toprule
JT05 / T\"auber target & CFAC status \\
\midrule
$\sim 10$ 1PI diagrams to enumerate         & \textbf{0} (Lagrange inversion delivers all $1+2+5 = 8$ topologies) \\
$\sim 10$ symmetry factors to assign        & \textbf{0} (Lagrange overcount automatic) \\
8--10 distinct loop integrals               & \textbf{3 master integrals} (IBP, two are textbook) \\
4 double-pole rationals (Eq.~57)            & \textbf{closed-form} via Hopf antipode \\
4 simple-pole rationals (Eq.~57)            & \textbf{12 CFAC rationals} via IBP closure \\
3 of 4 RG functions (Eq.~58a/b/c)           & \textbf{derived} via T\"auber relation, zero residual \\
4th RG function $\beta(u)$ at 2-loop (Eq.~58d) & 1-loop derived; 2-loop coefficient pending MSbar pole-closure (bookkeeping) \\
Fixed point $u^*$ (Eq.~59)                  & \textbf{derived}, zero residual \\
4 critical exponents (Eq.~60)               & \textbf{derived}, zero residual \\
3 master integral \emph{values}             & input (this is where AC stops and \emph{integration} begins) \\
\bottomrule
\end{tabular}
\end{center}

\section*{5.\ Caveat: AC counts, but doesn't autonomously select 1PI sectors}

The pipeline as we run it depends on QFT inputs we hand to AC: which external-leg sector is being renormalised (2-point for $Z_\psi$, 3-point for $Z_u$, etc.), which loop order, and the 1PI restriction. Lagrange inversion gives all $\phi$-trees at order $z^n$; the assignment of those counts to specific Z-factors is presently a hand-mapping in \texttt{gribov.tex \S 3.3} (e.g.\ ``$[z^7] = 5 = 3$ 1PI vertex topologies $+\,2$ reducible trees''). So in the current pipeline the relevance filter is QFT, not AC.

\section*{6.\ Proposed: combinatorial selection of relevant graphs}

\emph{This section sketches what would mechanise the topology-to-Z-factor mapping. We have not implemented it; sending it to you because you may want to push on it.}

\medskip
The QFT relevance filter (external legs, loop order, 1PI) admits a fully combinatorial encoding via multivariate marking. Build, in place of the single-variable $G(z)$, a multivariate generating function
\[
\mathcal{G}(z;\,\xi_\psi, \xi_{\tilde\psi};\,\alpha) \;=\; \sum_{n,e_\psi,e_{\tilde\psi},s} c_{n,e_\psi,e_{\tilde\psi},s}\;
z^{n}\,\xi_\psi^{e_\psi}\,\xi_{\tilde\psi}^{e_{\tilde\psi}}\,\alpha^{s},
\]
with $z$ counting vertex insertions (loop order), $\xi_\psi, \xi_{\tilde\psi}$ counting external legs of each type, and $\alpha$ tracking rapidity sign (already done in CFAC C2 via $G(z,\alpha)=z(1+\alpha G^2)$ at $\alpha = -1$). Then each Z-factor in JT05 is a \emph{specific bidegree projection} of $\mathcal{G}$:
\begin{align*}
\text{1-loop self-energy} &\;=\; [z^1\,\xi_\psi\,\xi_{\tilde\psi}]\,\mathcal{G}_{\text{1PI}}(z;\xi,\alpha{=}{-}1),\\
\text{1-loop vertex} &\;=\; [z^1\,\xi_\psi\,\xi_{\tilde\psi}^2]\,\mathcal{G}_{\text{1PI}}(z;\xi,\alpha{=}{-}1),\\
\text{2-loop }\Sigma &\;=\; [z^2\,\xi_\psi\,\xi_{\tilde\psi}]\,\mathcal{G}_{\text{1PI}}(z;\xi,\alpha{=}{-}1), \quad\text{etc.}
\end{align*}

The 1PI restriction is itself combinatorial: $\mathcal{G}_{\text{1PI}}$ is the Legendre transform of the connected $\mathcal{W} = \ln\mathcal{Z}$, and Legendre transform on multivariate generating functions is a closed-form algebraic operation (Flajolet--Sedgewick, multivariate AC). So once $\mathcal{Z}$ is built from $\phi$, every Z-factor is a $\pi_X$ projection on $\Gamma = \text{Legendre}(\ln\mathcal{Z})$.

The power-counting filter (which sectors are UV-divergent at $d_c = 4{-}\varepsilon$) is also combinatorial: each monomial of $\mathcal{G}$ has a fixed superficial divergence $\omega = dL - 2I + n_\partial$, computable from its $(z, \xi)$-degree pattern alone. Filter to $\omega \ge 0$ and the divergent sectors fall out.

\paragraph{What this would buy.} An autonomous combinatorial selector:
\[
\big(\text{action}\big) \;\to\; \big(\phi,\,\text{leg dictionary}\big) \;\to\; \mathcal{G}(z;\xi;\alpha)
\;\to\; \pi_X[\Gamma_{\text{1PI}}]\;\xrightarrow{\text{filter }\omega\ge 0}\; \{\text{Z-factors}_X\}
\]
with no human-mapped topology table.

\paragraph{What it would cost.} The Legendre transform on a multivariate generating function is the main piece of unwritten code. The marker $\xi_\psi, \xi_{\tilde\psi}$ extension of $G(z,\alpha)$ is a dozen lines of sympy. Power-counting filter is one inequality on monomial degrees. Estimated effort: a focused day of work to handle 2-loop DP, plus checks against the existing topology table.

\paragraph{Why we haven't done it yet.} For 2-loop DP the topology list is $\sim 5$ graphs --- the hand-mapping is faster than building the autonomous selector. For higher loops or a different theory (your H-vertex model, Manna, BARW), the autonomous selector pays off. Worth doing if/when one of those becomes the active target.

\section*{7.\ $\varphi^4$ in Doi--Peliti}

The same template applies: replace $\phi(G) = 1+G^2$ with $\phi(G) = 1+G^4$. Lagrange inversion, Hopf antipode, IBP closure, T\"auber relation, sympy series solution --- everything except the master integral values is coupling-agnostic. We have not run this case end-to-end; happy to do so on receipt of your preferred action / vertex normalisation.

\section*{8.\ Where the work lives}

\begin{itemize}[leftmargin=2em,itemsep=2pt]
  \item \texttt{paper/cfac/gribov.tex} --- the full write-up (CFAC theorem, IBP reduction, Hopf antipode, assembly).
  \item \texttt{rdft/ac/gribov/} --- mechanised pipeline:
    \texttt{assembly.py} (Lagrange + algebra),
    \texttt{actrick.py} (Hopf antipode),
    \texttt{ibp\_coefficients.py} (12 IBP rationals),
    \texttt{two\_loop.py} (full pipeline to Eq.~60).
  \item \texttt{rdft/ac/gribov/run\_all\_tests.py} --- 10 tests, all pass, including the JT05 Eq.~(60) match.
  \item \texttt{letters/.../eq58\_from\_z.py} (this letter directory) --- derives Eq.~(58a/b/c) from the CFAC $Z$-factors via the T\"auber relation.
\end{itemize}

\medskip
\noindent The companion note Q1 deals with your H-vertex ladder; the connection is that the same template (vertex polynomial $\to$ boundary-projected sector $\to$ closed-form residue / eigenvalue) handles both observables. Counts on one end, exponents on the other.

\end{document}
